\scanpage{74}
Les formules (10) (11) se réduisent donc à la formule
\[
(18)\qquad g_{zy}(u)g_{yx}(u)=g_{zx}(u)\qquad x,y,z\in T_K
\]
\margin{On a laissé tomber l'exposant $T$ : $g_{yx}^{T/X}$, qui est \unsure{nécessaire} surtout si $T\subset X\ldots$}
et la formule (12) à la formule itérée de celle-ci
\[
(19)\qquad g_{x_n,x_0}=g_{x_n,x_{n-1}}\,g_{x_{n-1},x_{\unsure{n-2}}}\cdots g_{x_1,x_0}
\]
--- quant aux formules de passage (9), elles \deleted{\unsure{tombent}} disparaissent !

Bien entendu, les invariants $g_l(u)$ \inserted{ou $f_l(u)$} peuvent être considérés comme décrits par des invariants précédents du type $g_{y,x}^T$, en prenant \deleted{\unsure{$T$}} p.ex. $T=[0,1]$, $p$ une application $[0,1]\longrightarrow X$ \unsure{donnant} naissance \deleted{au} chemin $l$ (plutôt, à la classe de chemins\ldots) de sorte qu'on aura
\[
(20)\qquad f_l(u)=\varphi_x\bigl(g_{y_0,x_0}^T(u)\bigr),\qquad g_l(u)=\varphi_y\bigl(g_{y_0,x_0}^T(u)\bigr)
\]
\deleted{$x_0=0\in T,\ y_0$}
\[
x_0=0,\ \deleted{\unsure{0}}\ y_0=1\qquad\in T=[0,1].
\]
\margin{Autre exemple ``\unsure{illuminant}'' : si $x_0\in X_K$ fixé, pour les $y\in X_K$ ``voisins'' de $\unsure{X|x_0}$, \unsure{choisissons} des chemins ``directs'' $x_0\longrightarrow y$ (``voisins'' de $x_0$) on trouve des invariants $g_{y,x_0}(u)\in\Pi_1(X,x_0)$.}

\deleted{La relation de $1$-cocycle}
\deleted{Bien entendu,}
\deleted{On peut}
Pour comparer les invariants $g_{x,y}^T$ associés à différents espaces $1$-connexes au dessus de $X$, il suffit de regarder le cas où on a un diag. comm.
\diagram{\begin{tikzcd}[ampersand replacement=\&,column sep=large]
T\arrow[rr,"f"]\arrow[dr,"p"'] \&\& T'\arrow[dl,"p'"]\\
\&X\&
\end{tikzcd}}
\[
(21)
\]
d'où $f_K:T_K\mathbin{-}T'_K$,
\margin{73}
